\chapter{Аналитическая часть}
%<пишем про тему и все, связанное с ней>

Задача коммивояжёра заключается в поиске минимального пути, при котором все узлы графа будут посещены~\cite{lit2}.
Гамильтонов цикл -- это модификация задачи коммивояжёра, добавляющая ещё одно условие к изначальной формулировке: начало обхода совпадет к его концом.

\section{Алгоритм полного перебора}

Суть алгоритма полного перебора лежит в его названии: все возможные пути графа последовательно перебираются, в результате находится либо самый оптимальный путь, либо его отсутствие~\cite{lit3}.

\section{Муравьиный алгоритм}

Идея муравьиного алгорима (Ant Colony Optimization, ACO) заключается в моделировании поведения колонии муравьёв, которые коллективно находят оптимальные пути, используя феромоны как механизм координации и обмена информацией~\cite{lit4}.

Каждые муравьиные сутки (итерация алгоритма) каждый из $m$ муравьёв помещается в случайную вершину $i_{k}$ и начинает движение. Следующая вершина маршрута $j$ выбирается в зависимости от вероятности~(формула \ref{pk}):
\begin{align}
	p_{ij}^{k} = \dfrac{\tau_{ij}^{\alpha} \cdot \eta_{ij}^{\beta}}{\sum_{l \in N_{i}^{k}}{\tau_{il}^{\alpha} \cdot \eta_{il}^{\beta}}}
	\label{pk}
\end{align}
где 
\begin{itemize}
	\item $\tau_{ij}$ -- концентрации феромона на дуге из вершины $i$ в  вершину $j$;
	\item $\eta_{ij} = \frac{1}{d_{ij}}$ -- эвристика (исследование);
	\item $\alpha, \beta$ -- коэффициенты влияния феромона и жадности соответственно;
	\item $N_{i}^{k}$ -- множество непосещённых соседей муравья $k$ из вершины $i$.
\end{itemize}
По истечении суток полученные маршруты $L_{k}$ сравниваются с результатами предыдущих суток, а концентрации феромонов обновляются по правилу~(формула \ref{pher}):
\begin{align}
	\tau_{ij} = (1 - \rho)\tau_{ij} + \sum_{k=1}^{m}{\Delta\tau_{ij}},
	\label{pher}
\end{align}
где $\rho$ -- коэффициент испарения, а 
\begin{equation}
	\Delta\tau_{ij} = 
		\begin{cases}
			\frac{Q}{L_{k}}, & \text{если муравей $k$ прошёл по дуге $ij$},\\
			0, & \text{иначе}.
		\end{cases}
\end{equation}
$Q$ -- максимальное количество феромона, которое вычисляется относительно длин рёбер рассматриваемого графа.

По истечении предельного количества суток возвращается итоговый результат -- оптимальный путь графа.

\section{Задание по варианту}

Ор-граф, без элитных муравьёв, вокруг света за 80 дней (на рёбрах затраты по времени; есть предел -- 80 дней в человеческом измерении, не совпадающем с муравьиным). Гамильтонов цикл.

% 80 * 24 * 60 * 60 = 6912000 c

\section*{Вывод}

В данном разделе мной был проведен анализ теоретических материалов, объясняющих принцип работы алгоритмов.

\clearpage
